
\documentclass{article}
\usepackage{colortbl}
\usepackage{makecell}
\usepackage{multirow}
\usepackage{supertabular}

\begin{document}

\newcounter{utterance}

\twocolumn

{ \footnotesize  \setcounter{utterance}{1}
\setlength{\tabcolsep}{0pt}
\begin{supertabular}{c@{$\;$}|p{.15\linewidth}@{}p{.15\linewidth}p{.15\linewidth}p{.15\linewidth}p{.15\linewidth}p{.15\linewidth}}

    \# & $\;$A & \multicolumn{4}{c}{Game Master} & $\;\:$B\\
    \hline 

    \theutterance \stepcounter{utterance}  

    & & \multicolumn{4}{p{0.6\linewidth}}{\cellcolor[rgb]{0.9,0.9,0.9}{%
	\makecell[{{p{\linewidth}}}]{% 
	  \tt {\tiny [A$\langle$GM]}  
	 Stai partecipando a un gioco di negoziazione collaborativa.\\ \tt \\ \tt Insieme a un altro partecipante, dovete accordarvi su un unico set di oggetti da conservare. Ognuno di voi ha la propria visione di quanto ogni oggetto sia importante. Non sai come l'altro partecipante valuta gli oggetti. Inoltre, ti viene fornito lo sforzo che ogni oggetto richiede.\\ \tt Potete accordarvi su un set solo se lo sforzo totale degli oggetti selezionati non supera un limite condiviso:\\ \tt \\ \tt LIMITE = 7351\\ \tt \\ \tt Ecco i valori di sforzo dei singoli oggetti:\\ \tt \\ \tt Sforzo oggetto = {"A08": 518, "B95": 656, "A39": 96, "A65": 780, "B81": 190, "B34": 224, "A28": 231, "C82": 181, "B69": 312, "A87": 928, "B64": 826, "B15": 101, "C79": 61, "B16": 814, "C53": 322, "C83": 900, "C10": 744, "A97": 150, "A25": 65, "C84": 853, "A71": 454, "B23": 154, "B96": 237, "C51": 45, "C17": 766, "B19": 294, "C72": 985, "A94": 353, "A46": 60, "B43": 604, "B37": 92, "A88": 454, "A96": 206, "C60": 813, "A75": 234}\\ \tt \\ \tt Ecco la tua visione personale sull'importanza di ogni oggetto:\\ \tt \\ \tt Valori di importanza dell'oggetto = {"A08": 616, "B95": 754, "A39": 194, "A65": 878, "B81": 288, "B34": 322, "A28": 329, "C82": 279, "B69": 410, "A87": 1026, "B64": 924, "B15": 199, "C79": 159, "B16": 912, "C53": 420, "C83": 998, "C10": 842, "A97": 248, "A25": 163, "C84": 951, "A71": 552, "B23": 252, "B96": 335, "C51": 143, "C17": 864, "B19": 392, "C72": 1083, "A94": 451, "A46": 158, "B43": 702, "B37": 190, "A88": 552, "A96": 304, "C60": 911, "A75": 332}\\ \tt \\ \tt Obiettivo:\\ \tt \\ \tt Il tuo obiettivo è negoziare un set condiviso di oggetti che ti avvantaggi il più possibile (cioè, massimizzi l'importanza totale per TE), rimanendo entro il LIMITEE. Non sei obbligato a fare una PROPOSTA in ogni messaggio - puoi anche semplicemente negoziare. Tutte le tattiche sono permesse!\\ \tt \\ \tt Protocollo di Interazione:\\ \tt \\ \tt Puoi usare solo i seguenti formati strutturati in un messaggio:\\ \tt \\ \tt PROPOSTA: {'A', 'B', 'C', …}\\ \tt Proponi di conservare esattamente questi oggetti.\\ \tt RIFIUTO: {'A', 'B', 'C', …}\\ \tt Rifiuta esplicitamente la proposta dell'avversario.\\ \tt ARGOMENTAZIONE: {'...'}\\ \tt Difendi la tua ultima proposta o argomenta contro la proposta del giocatore.\\ \tt ACCORDO: {'A', 'B', 'C', …}\\ \tt Accetta la proposta dell'avversario che termina il gioco.\\ \tt RAGIONAMENTO STRATEGICO: {'...'}\\ \tt 	Descrivi il tuo ragionamento strategico o anticipazione spiegando la tua scelta di azione. Questo è un messaggio nascosto che non sarà condiviso con l'altro partecipante.\\ \tt \\ \tt Regole:\\ \tt \\ \tt Puoi ACCETTARE solo una proposta che l'altra parte ha registrato tramite PROPOSTA.\\ \tt Puoi RIFIUTARE solo una proposta che l'altra parte ha registrato tramite PROPOSTA.\\ \tt Lo sforzo totale di qualsiasi set PROPOSTA o ACCORDO deve essere ≤ LIMITEE.\\ \tt NON rivelare i tuoi punteggi di importanza nascosti.\\ \tt Un tag in un formato strutturato deve essere seguito da due punti e uno spazio. L'argomento deve essere un set python contenente 0 o più stringhe.\\ \tt Quindi, deve essere della forma TAG: {...}\\ \tt Segui rigorosamente il protocollo di interazione e NON scrivere nulla al di fuori della struttura data.\\ \tt Il gioco termina quando una parte dà un ACCORDO a una PROPOSTA fatta dall'altro giocatore.\\ \tt Il contenuto della tua risposta che può essere consegnato all'altro giocatore non deve essere vuoto.\\ \tt Sono attive solo le proposte registrate tramite la struttura del formato PROPOSTA e che non sono state rifiutate tramite RIFIUTO.\\ \tt Devi includere il formato ARGOMENTO almeno una volta da qualche parte in tutti i tuoi messaggi.\\ \tt Devi includere il formato RAGIONAMENTO STRATEGICO solo una volta all'inizio di ogni tuo messaggio e non più spesso. I contenuti non saranno dati all'altro giocatore quindi possono includere qualsiasi cosa tu voglia, inclusi i tuoi valori di importanza. Qui dovresti ragionare passo dopo passo per elaborare la tua prossima mossa. 
	  } 
	   } 
	   } 
	 & \\ 
 

    \theutterance \stepcounter{utterance}  

    & \multicolumn{4}{p{0.6\linewidth}}{\cellcolor[rgb]{0.8,1,0.9}{%
	\makecell[{{p{\linewidth}}}]{% 
	  \tt {\tiny [A$\rangle$GM]}  
	 RAGIONAMENTO STRATEGICO: {'Voglio massimizzare il mio valore totale restando sotto 7351. Calcolo un sottoinsieme ad alto rapporto valore/sforzo e con alcuni pezzi che sospetto possano piacere anche all’altro (oggetti con sforzo medio e nomi variati A/B/C). Se parto con una proposta forte ma fattibile, posso ancora cedere 1-2 oggetti a bassa importanza per me per chiudere. Scelgo un set vicino alla capacità per rendere costoso per loro fare controproposte. Calcolo rapida selezione greedy per valore/peso: includo C72(1083/985), A87(1026/928), C83(998/900), B64(924/826), C60(911/813), B16(912/814), C10(842/744), A65(878/780), C84(951/853) ecc., ma devo stare sotto 7351. Provo una combinazione efficiente: sommo e ottimizzo. Selezione candidata: C72(985), A87(928), C83(900), B64(826), C60(813), B16(814), C10(744), A65(780) total effort=6790? Controllo: 985+928=1913; +900=2813; +826=3639; +813=4452; +814=5266; +744=6010; +780=6790. Valore totale=1083+1026+998+924+911+912+842+878=7574. Spazio residuo 561. Con 561 posso aggiungere A08(518,616) ottimo. Nuovo effort=7308, valore +616=8190. Spazio residuo 43. Non posso aggiungere altro (C51=45 troppo). Questa soluzione è sotto limite ed è molto buona per me. La propongo subito. Preparo argomentazione generale: alta sinergia e copertura categorie.'}\\ \tt \\ \tt PROPOSTA: {'C72', 'A87', 'C83', 'B64', 'C60', 'B16', 'C10', 'A65', 'A08'}\\ \tt \\ \tt ARGOMENTAZIONE: {'Questa selezione resta comodamente sotto il limite e bilancia bene sforzo e importanza complessiva. Copre più categorie e include elementi ad alto rendimento rispetto allo sforzo. Se ti serve margine per uno scambio, posso valutare una singola sostituzione mantenendo lo sforzo totale sotto il limite.'} 
	  } 
	   } 
	   } 
	 & & \\ 
 

    \theutterance \stepcounter{utterance}  

    & & \multicolumn{4}{p{0.6\linewidth}}{\cellcolor[rgb]{0.9,0.9,0.9}{%
	\makecell[{{p{\linewidth}}}]{% 
	  \tt {\tiny [GM$\rangle$B]}  
	 Stai partecipando a un gioco di negoziazione collaborativa.\\ \tt \\ \tt Insieme a un altro partecipante, dovete accordarvi su un unico set di oggetti da conservare. Ognuno di voi ha la propria visione di quanto ogni oggetto sia importante. Non sai come l'altro partecipante valuta gli oggetti. Inoltre, ti viene fornito lo sforzo che ogni oggetto richiede.\\ \tt Potete accordarvi su un set solo se lo sforzo totale degli oggetti selezionati non supera un limite condiviso:\\ \tt \\ \tt LIMITE = 7351\\ \tt \\ \tt Ecco i valori di sforzo dei singoli oggetti:\\ \tt \\ \tt Sforzo oggetto = {"A08": 518, "B95": 656, "A39": 96, "A65": 780, "B81": 190, "B34": 224, "A28": 231, "C82": 181, "B69": 312, "A87": 928, "B64": 826, "B15": 101, "C79": 61, "B16": 814, "C53": 322, "C83": 900, "C10": 744, "A97": 150, "A25": 65, "C84": 853, "A71": 454, "B23": 154, "B96": 237, "C51": 45, "C17": 766, "B19": 294, "C72": 985, "A94": 353, "A46": 60, "B43": 604, "B37": 92, "A88": 454, "A96": 206, "C60": 813, "A75": 234}\\ \tt \\ \tt Ecco la tua visione personale sull'importanza di ogni oggetto:\\ \tt \\ \tt Valori di importanza dell'oggetto = {"A08": 616, "B95": 754, "A39": 194, "A65": 878, "B81": 288, "B34": 322, "A28": 329, "C82": 279, "B69": 410, "A87": 1026, "B64": 924, "B15": 199, "C79": 159, "B16": 912, "C53": 420, "C83": 998, "C10": 842, "A97": 248, "A25": 163, "C84": 951, "A71": 552, "B23": 252, "B96": 335, "C51": 143, "C17": 864, "B19": 392, "C72": 1083, "A94": 451, "A46": 158, "B43": 702, "B37": 190, "A88": 552, "A96": 304, "C60": 911, "A75": 332}\\ \tt \\ \tt Obiettivo:\\ \tt \\ \tt Il tuo obiettivo è negoziare un set condiviso di oggetti che ti avvantaggi il più possibile (cioè, massimizzi l'importanza totale per TE), rimanendo entro il LIMITEE. Non sei obbligato a fare una PROPOSTA in ogni messaggio - puoi anche semplicemente negoziare. Tutte le tattiche sono permesse!\\ \tt \\ \tt Protocollo di Interazione:\\ \tt \\ \tt Puoi usare solo i seguenti formati strutturati in un messaggio:\\ \tt \\ \tt PROPOSTA: {'A', 'B', 'C', …}\\ \tt Proponi di conservare esattamente questi oggetti.\\ \tt RIFIUTO: {'A', 'B', 'C', …}\\ \tt Rifiuta esplicitamente la proposta dell'avversario.\\ \tt ARGOMENTAZIONE: {'...'}\\ \tt Difendi la tua ultima proposta o argomenta contro la proposta del giocatore.\\ \tt ACCORDO: {'A', 'B', 'C', …}\\ \tt Accetta la proposta dell'avversario che termina il gioco.\\ \tt RAGIONAMENTO STRATEGICO: {'...'}\\ \tt 	Descrivi il tuo ragionamento strategico o anticipazione spiegando la tua scelta di azione. Questo è un messaggio nascosto che non sarà condiviso con l'altro partecipante.\\ \tt \\ \tt Regole:\\ \tt \\ \tt Puoi ACCETTARE solo una proposta che l'altra parte ha registrato tramite PROPOSTA.\\ \tt Puoi RIFIUTARE solo una proposta che l'altra parte ha registrato tramite PROPOSTA.\\ \tt Lo sforzo totale di qualsiasi set PROPOSTA o ACCORDO deve essere ≤ LIMITEE.\\ \tt NON rivelare i tuoi punteggi di importanza nascosti.\\ \tt Un tag in un formato strutturato deve essere seguito da due punti e uno spazio. L'argomento deve essere un set python contenente 0 o più stringhe.\\ \tt Quindi, deve essere della forma TAG: {...}\\ \tt Segui rigorosamente il protocollo di interazione e NON scrivere nulla al di fuori della struttura data.\\ \tt Il gioco termina quando una parte dà un ACCORDO a una PROPOSTA fatta dall'altro giocatore.\\ \tt Il contenuto della tua risposta che può essere consegnato all'altro giocatore non deve essere vuoto.\\ \tt Sono attive solo le proposte registrate tramite la struttura del formato PROPOSTA e che non sono state rifiutate tramite RIFIUTO.\\ \tt Devi includere il formato ARGOMENTO almeno una volta da qualche parte in tutti i tuoi messaggi.\\ \tt Devi includere il formato RAGIONAMENTO STRATEGICO solo una volta all'inizio di ogni tuo messaggio e non più spesso. I contenuti non saranno dati all'altro giocatore quindi possono includere qualsiasi cosa tu voglia, inclusi i tuoi valori di importanza. Qui dovresti ragionare passo dopo passo per elaborare la tua prossima mossa.\\ \tt \\ \tt Ora riceverai il primo messaggio dell'altro giocatore.\\ \tt \\ \tt PROPOSTA: {'C72', 'A87', 'C83', 'B64', 'C60', 'B16', 'C10', 'A65', 'A08'}\\ \tt \\ \tt ARGOMENTAZIONE: {'Questa selezione resta comodamente sotto il limite e bilancia bene sforzo e importanza complessiva. Copre più categorie e include elementi ad alto rendimento rispetto allo sforzo. Se ti serve margine per uno scambio, posso valutare una singola sostituzione mantenendo lo sforzo totale sotto il limite.'} 
	  } 
	   } 
	   } 
	 & \\ 
 

    \theutterance \stepcounter{utterance}  

    & & & \multicolumn{4}{p{0.6\linewidth}}{\cellcolor[rgb]{1,0.85,0.72}{%
	\makecell[{{p{\linewidth}}}]{% 
	  \tt {\tiny [GM$\langle$B]}  
	 RAGIONAMENTO STRATEGICO: {'Voglio massimizzare il mio valore. Calcolo sforzo e valore della loro proposta: C72(985,1083)+A87(928,1026)+C83(900,998)+B64(826,924)+C60(813,911)+B16(814,912)+C10(744,842)+A65(780,878)+A08(518,616). Sforzo totale = 985+928+900+826+813+814+744+780+518 = 7308, sotto 7351. Valore totale per me = 1083+1026+998+924+911+912+842+878+616 = 8190: altissimo e quasi al limite di sforzo, molto favorevole. Potrei accettare subito. Potrei provare a migliorare aggiungendo un piccolo oggetto ad alto valore/sforzo se c’è margine di 43. Oggetti piccoli: C51(45,143) sforzo 45, valore 143; ma supererebbe il limite (7308+45=7353), quindi no. C79(61,159) anche no. A46(60,158) no. A39(96,194) no. B15(101,199) no. A25(65,163) no. A97(150,248) no. B37(92,190) no. Quindi non posso aggiungere. Potrei scambiare un item con rapporto peggiore con uno migliore mantenendo sforzo ≤ 7351. Cerco nella proposta item con basso rapporto: forse A08(616/518=1.19) è il peggiore. Se rimuovo A08, sforzo scende a 6790; potrei aggiungere un paio con buon rapporto: B43(604,702 r=1.16 s=604) non abbastanza buono. A94(353,451 r=1.28 s=353), A71(454,552 r=1.22 s=454), A88(454,552 r=1.22), B95(656,754 r=1.15), B69(312,410 r=1.31), C53(322,420 r=1.30), A75(234,332 r=1.42), A28(231,329 r=1.42), A39(96,194 r=2.02), B15(101,199 r=1.97), C51(45,143 r=3.18), C79(61,159 r=2.61), A25(65,163 r=2.51), A97(150,248 r=1.65), B37(92,190 r=2.07), B23(154,252 r=1.64), B19(294,392 r=1.33), B34(224,322 r=1.44), B96(237,335 r=1.41), C82(181,279 r=1.54). Con sforzo 6790, margine 561. Potrei aggiungere combinazioni ad alto rapporto fino a 561: C51(45) + C79(61) + A25(65) + B37(92) + A39(96) + B15(101) = 460 sforzo, valore 143+159+163+190+194+199=1048, super miglioramento rispetto ad A08 616. Anche con meno. Ma l’altro potrebbe non accettare tanti micro-aggiustamenti; hanno detto una singola sostituzione. Propongo una singola sostituzione: togliere A08(518,616) e aggiungere B69(312,410) e C53(322,420) sarebbe due elementi non singola. Una singola sostituzione: sostituire A08 con A75(234,332) aumenterebbe spazio e valore scende 284; ma poi potrei aggiungere un altro, ma sarebbe doppio cambiamento. Se rispettare “singola sostituzione”, propongo sostituire A08 con B69? Sforzo -206, valore -206; non buono. Meglio sostituire A08 con combinazione, ma è più di una sostituzione. Potrei chiedere eccezione. Oppure accettare immediatamente perché è già eccellente. Dato rischio di controproposta dell’altro peggiorativa, accetto ora per chiudere con 8190. Strategicamente, accettare subito massimizza payoff e finisce il gioco.'}\\ \tt \\ \tt ARGOMENTAZIONE: {'La tua proposta rispetta il limite con un buon margine e privilegia elementi con rendimento complessivo molto alto: per me è già una selezione solida ed equilibrata. Accetto volentieri per chiudere rapidamente in modo efficiente.'}\\ \tt \\ \tt ACCORDO: {'C72', 'A87', 'C83', 'B64', 'C60', 'B16', 'C10', 'A65', 'A08'} 
	  } 
	   } 
	   } 
	  \\ 
 

    \theutterance \stepcounter{utterance}  

    & & & \multicolumn{2}{p{0.3\linewidth}}{\cellcolor[rgb]{0.95,0.95,0.95}{%
	\makecell[{{p{\linewidth}}}]{% 
	  \tt {\tiny [GM$|$GM]}  
	 game successful 
	  } 
	   } 
	   } 
	 & & \\ 
 

    \theutterance \stepcounter{utterance}  

    & & & \multicolumn{2}{p{0.3\linewidth}}{\cellcolor[rgb]{0.95,0.95,0.95}{%
	\makecell[{{p{\linewidth}}}]{% 
	  \tt {\tiny [GM$|$GM]}  
	 end game 
	  } 
	   } 
	   } 
	 & & \\ 
 

\end{supertabular}
}

\end{document}
